% mycsrf 'for beeing included' snippet template
%
% (c) Karsten Reincke, Frankfurt a.M. 2012, ff.
%
% This text is licensed under the Creative Commons Attribution 3.0 Germany
% License (http://creativecommons.org/licenses/by/3.0/de/): Feel free to share
% (to copy, distribute and transmit) or to remix (to adapt) it, if you respect
% how you must attribute the work in the manner specified by the author(s):
% \newline
% In an internet based reuse please link the reused parts to mycsrf.fodina.de
% and mention the original author Karsten Reincke in a suitable manner. In a
% paper-like reuse please insert a short hint to mycsrf.fodina.de and to the
% original author, Karsten Reincke, into your preface. For normal quotations
% please use the scientific standard to cite
%

\section{Vom Frontend nach \LaTeX: die Toolketten}

Wenn wir nun einen Schritt zurücktreten und uns auf die
Tools konzentrieren, denen wir drei oder mehr Sterne gegeben haben, dann
gewinnen wir folgende Übersicht:


\begin{tikzpicture}
\footnotesize

\tikzstyle{every node}=[trapezium, draw, minimum width=1cm,trapezium left angle=70, trapezium right angle=110]

% BOTTOM UP DESCRIPTION of the component IO structure:
% THE FINALLY GENERATED PDF OUTPOUT ->
\file{\pdfBgColor}{black}{PDF}{5.5}{0}{.pdf}

% -> CREATED BY THE LATEX IDEs ->
\backend{$\star\star\star$}{\abcColor}{LTXABC}{1}{2}{\LaTeX\,\&\,ABC}
\backend{$\star\star\star\star$}{\mtexColor}{LTXMTEX}{5.5}{2}{\LaTeX\,\&\,MusiX\TeX}
\backend{$\star\star\star\star\star$}{\lyColor}{LTXLIPO}{10.4}{2}{\LaTeX\,\&\,LilyPond}

% -> which take the CONVERTER OUTPUT ->
\file{\abcBgColor}{\abcColor}{CO-ABC}{1}{3}{.abc}
\file{\mtexBgColor}{\mtexColor}{CO-MTEX}{5.5}{3}{.mtex}
\file{\lyBgColor}{Blue}{CO-LY}{10.4}{3}{.ly}

% written by the CONVERTERs.
\converter{\mxmlInColor}{\abcOutColor}{XML2ABC}{0}{7}{xml2abc}
\converter{\mxmlInColor}{\abcOutColor}{MXML2ABC}{1.5}{6}{mxml2abc}
\converter{\mxmlInColor}{\pmxOutColor}{MXML2PMX}{3.6}{8}{mxml2pmx}
\converter{\mxmlInColor}{\lyOutColor}{XML2LY}{4}{11}{xml2ly}
\converter{\mxmlInColor}{\lyOutColor}{MXML2LY}{4.5}{14}{mxml2ly}

\converter{\abcInColor}{\mtexOutColor}{ABC2MTEX}{3}{4}{abc2mtex}
\converter{\abcInColor}{\lyOutColor}{ABC2LY}{5}{12}{abc2ly}

\converter{\lyInColor}{\abcOutColor}{LY2ABC}{5.8}{8}{ly2abc}
\converter{\pmxInColor}{\mtexOutColor}{PMXAB}{5.5}{4}{pmxab}


% These converters takte the FRONTEND OUTPUT ->

\file{\abcBgColor}{\abcColor}{FO-ABC}{0}{15}{.abc}
\file{\mxmlBgColor}{\mxmlColor}{FO-MXML}{1.5}{15}{.mxml}
\file{\pmxBgColor}{Brown}{FO-PMX}{5.5}{5.6}{.pmx}
\file{\lyBgColor}{Blue}{FO-LY}{7.2}{11}{.ly}

% created by the FRONTENDS ->

\frontend{\abcBgColor}{Black}{EASYABC}{0}{17.5}{EasyABC}{$\star\star\star\star$};
\frontend{\lyBgColor}{Black}{ROSEGARDEN}{2.2}{18.5}{Rosegarden}{$\star\star\star$};
\frontend{\multiBgColor}{Black}{MUSESCORE}{4.5}{18.5}{MuseScore}{$\star\star\star\star\star$};
\frontend{\lyBgColor}{Black}{DENEMO}{6.5}{18.5}{Denemo}{$\star\star\star$};
\frontend{\lyBgColor}{Black}{CANORUS}{8.3}{18.5}{Canorus}{$\star\star\star$};

\frontend{\multiBgColor}{Black}{EDITOR}{7.7}{7.2}{Text-Editor}{$\star$};
\frontend{\lyBgColor}{Black}{ELYSIUM}{9.8}{8}{Elysium}{$\star\star\star\star$};
\frontend{\lyBgColor}{Black}{FRESCOBALDI}{11.9}{8}{Frescobaldi}{$\star\star\star\star\star$};

% which are used by the Composer:

\node[circle,draw=none] (Composer) at (6.5,21) 
{ \includegraphics{./logos/decider-220x320.eps}};
  
% TOP DOWN linkings of the component IO structure:  
  
% (1) The composer uses the frontends.
\draw[->] (Composer) to [out=180,in=90] (EASYABC);  
\draw[->] (Composer) to [out=180,in=90] (ROSEGARDEN); 
\draw[->] (Composer) to [out=270,in=90] (MUSESCORE);  
\draw[->] (Composer) to [out=270,in=90] (DENEMO);  
\draw[->] (Composer) to [out=270,in=90] (CANORUS);  
\draw[->] (Composer) to [out=0,in=90] (ELYSIUM);  
\draw[->] (Composer) to [out=0,in=90] (FRESCOBALDI);  
\draw[->] (Composer) to [out=0,in=90] (EDITOR);  

 
% (2.A) The frontends write the frontend output  
\draw[->,\abcColor] (EASYABC) to [out=270,in=90] (FO-ABC); 
\draw[->,\abcColor] (EASYABC) to [out=240,in=140] (CO-ABC); 
\draw[->,\mxmlColor] (EASYABC) to [out=300,in=90] (FO-MXML); 

\draw[->,\lyColor] (ROSEGARDEN) to [out=300,in=90] (FO-LY); 
\draw[->,\mxmlColor] (ROSEGARDEN) to [out=240,in=90] (FO-MXML);

\draw[->,\mxmlColor] (MUSESCORE) to [out=240,in=90] (FO-MXML); 

\draw[->,\mxmlColor] (CANORUS) to [out=240,in=90] (FO-MXML); 
\draw[->,\lyColor] (CANORUS) to [out=270,in=90] (FO-LY); 

\draw[->,\mxmlColor] (DENEMO) to [out=240,in=90] (FO-MXML); 
\draw[->,\lyColor] (DENEMO) to [out=270,in=90] (FO-LY); 


\draw[->,\abcColor] (EDITOR) to [out=270,in=80] (CO-ABC);
\draw[->,\mtexColor] (EDITOR) to [out=270,in=0] (CO-MTEX);
\draw[->,\pmxColor] (EDITOR) to [out=270,in=90] (FO-PMX);

\draw[->,\lyColor] (EDITOR) to [out=270,in=90] (CO-LY);
\draw[->,\lyColor] (ELYSIUM) to [out=270,in=90] (CO-LY); 
\draw[->,\lyColor] (FRESCOBALDI) to [out=270,in=90] (CO-LY); 

% (3) The frontend output becomes input of the converter

\draw[->,\abcColor] (FO-ABC) to [out=270,in=90] (ABC2MTEX);
\draw[->,\abcColor] (FO-ABC) to [out=270,in=90] (ABC2LY);

\draw[->,\mxmlColor] (FO-MXML) to [out=270,in=90] (XML2ABC);
\draw[->,\mxmlColor] (FO-MXML) to [out=270,in=90] (MXML2ABC);
\draw[->,\mxmlColor] (FO-MXML) to [out=270,in=90] (XML2LY);
\draw[->,\mxmlColor] (FO-MXML) to [out=270,in=90] (MXML2LY);
\draw[->,\mxmlColor] (FO-MXML) to [out=270,in=90] (MXML2PMX);

\draw[->,\lyColor] (FO-LY) to [out=270,in=90] (LY2ABC); 
%\draw[->,\lyColor] (FO-LY) to [out=270,in=90] (EDITOR); 
\draw[->,\lyColor] (FO-LY) to [out=270,in=90] (ELYSIUM); 
\draw[->,\lyColor] (FO-LY) to [out=270,in=90] (FRESCOBALDI); 

% (3) The converters (re)write their output
% (3.A) as input for the textual frontends
\draw[->,\lyColor] (ABC2LY) to [out=270,in=90] (FO-LY);
\draw[->,\lyColor] (MXML2LY) to [out=270,in=90] (FO-LY); 
\draw[->,\lyColor] (XML2LY) to [out=270,in=90] (FO-LY); 

% (3.B) as input for the backends

\draw[->,\abcColor] (LY2ABC) to [out=270,in=90] (CO-ABC);
\draw[->,\abcColor] (XML2ABC) to [out=270,in=90] (CO-ABC);
\draw[->,\abcColor] (MXML2ABC) to [out=270,in=90] (CO-ABC);

\draw[->,\pmxColor] (MXML2PMX) to [out=270,in=90] (FO-PMX);
\draw[->,\mtexColor] (ABC2MTEX) to [out=270,in=180] (CO-MTEX);

% special case MusiXTeX
\draw[->,\pmxColor] (FO-PMX) to [out=270,in=90] (PMXAB);
\draw[->,\mtexColor] (PMXAB) to [out=270,in=90] (CO-MTEX);

% 4) The backends take the (converter) output as input
\draw[->,\abcColor] (CO-ABC) to [out=270,in=90] (LTXABC);
\draw[->,\mtexColor] (CO-MTEX) to [out=270,in=90] (LTXMTEX);
\draw[->,\lyColor] (CO-LY) to [out=270,in=90] (LTXLIPO);

% 5) end write the final PDF

\draw[->,Black] (LTXABC) to [out=270,in=90] (PDF);
\draw[->,Black] (LTXMTEX) to [out=270,in=90] (PDF);
\draw[->,Black] (LTXLIPO) to [out=270,in=90] (PDF);
\end{tikzpicture}

Wir kennen demnach vier echte graphische Frontends -- \acc{Rosegarden},
\acc{MuseScore}, \acc{Denemo} und \acc{Canorus} --, drei semi-graphische
Frontends -- \acc{EasyABC}, \acc{Elysium} und \acc{Frescobaldi} --, sowie eine
Reihe von Konvertern. Unglücklicherweise haben wir lernen müssen, dass es keinen
graphischen oder semi-graphischen Editor für das \acc{MusiX\TeX}-System gibt.
Deshalb müssen wir sogar den reinen Texteditor berücksichtigen. Außerdem haben
wir festgestellt, dass keines der echten graphischen Frontends eine optimale
Dokumentation der Harmonieanalyse bietet. Und schließlich hatte sich schon
vorher ergeben, dass das \LaTeX-ABC-Backend nur die Integration begrenzt
komplexer Analysesymbole anbietet, während die Kombinationen von \LaTeX\ und
MusiX\TeX\ bzw. \LaTeX und LilyPond jeweils gut mit strukturell adäquaten
Analysesymbolen umgehen können. Allerdings sind diese Kennzeichnungen bei der
\LaTeX-MusiX\TeX-Variante aktuell visuell noch ansprechender, als bei der
\LaTeX-LilyPond-Variante.

Unsere 'Landkarte' ist noch eine bloß theoretische Skizze: die Wege ergaben sich
rein aus dem 'Versprechen' der Tools, bestimmte Formate zu schreiben und zu
lesen. Über die Praxis sagt das nichts. Die werden wir Weg für Weg evaluieren
müssen. Allerdings sollten wir die Güte eines Weges dann auch systematisch
vergleichbar in einer Zahl repräsentieren können. Glücklicherweise bietet sich
dafür unmittelbar ein Bewertungsverfahren an:

\begin{enumerate}
  \item Wir haben die Frontends und die Backends bereits bereits gemäß ihrer
  immanenten Fähigkeiten klassifiziert. Und jeder Weg beginnt mit einem Frontend
  und endet an einem Backend. Also sollte die Summe der Frontend- und
  Backend-Sterne etwas über die Qualität des Weges ausssagen.
  \item Allerdings ist eine erhöhte Komplexität des Weges ein Nachteil: je mehr
  Konvertierungsschritte nötig sind, desto größer ist der Aufwand für den
  Musikwissenschaftler. Es ist also angemessen, wenn wir von der Summe der
  Sterne die Anzahl der Konvertierungsschritte wieder abzuziehen: je direkter
  der Weg, desto besser.\footnote{Da aber zwischen einem Frontend und einem
  Backend immer mindestens ein Konvertierungsschritt liegt, der sich in der
  geschriebenen Outputdatei des Frontends manifestiert, die zum Input des
  Backends wird, brauchen wir letztlich von Summer der Sterne nur die jeweiligen
  Konvertierungsschritte - 1 zu subtrahieren.}
  \item In den Bewertung der Alternativen muss außerdem die Qualität des
  Ergebnisses (exzellent, akzeptabel, grenzwertig, inakzeptabel) eingerechnet
  werden.
  \item Ferner wird nicht nur die Anzahl der Konvertierungsschritte die
  (gefühlte) Nutzbarkeit eines Weges beeinflusses, sondern auch die technische
  Umsetzbarkeit (problemlos, stolpernd, beschränkt, unmöglich): Was vom Format
  her -- sozusagen auf dem Papier -- problemlos verknüpfbar ist, kann
  immer noch in technische Probleme laufen, die ungünstigenfalls die Nutzung
  gänzlich unmöglich machen.
\end{enumerate}
 
Die Qualität des Ergebnisses (\ra 3) und die technische Ausführbarkeit (\ra 4)
können in einen Faktor zusammengerechnet werden, mit dem das 'Sterne-Ergebnis'
(\ra 1 - \ra 2) gewichtet wird. Dieser Faktor ergibt sich geradezu natürlich:
Die Qualität eines Weges, der problemlos ein exzellentes Ergebnis erzeugt, wird
mit 1 multipliziert: er ist ohne Abstriche so gut, wie die Qualität der
beteiligten Frontends abzgl. der Komplexität. Wege, die inakzeptable Ergebnisse
bringen oder nicht umsetzbar sind, werden mit 0 gewichtet: sie sind nutzlos.
Und dazwischen entstehen folgende Abstufungen:

\begin{center}
\renewcommand{\arraystretch}{1.5}
\begin{tabulary}{14cm}{|C|R||C|C|C|C||}
\cline{3-6}
 \multicolumn{2}{c||}{\ } & \multicolumn{4}{c||}{Technische Ausführbarkeit} \\
\cline{3-6}
 \multicolumn{2}{c||}{\ } & problemlos & stolpernd & beschränkt & unmöglich \\
\hline
\multirow{4}{*}{\rotatebox{90}{Ergebnis}} & exzel\-lent & \centering{1} & 0.8 & 0.6 & 0 \\
\cline{2-6}
 & akzep\-tabel & 0.8 & 0.6 & 0.4 & 0 \\
\cline{2-6}
 & grenz\-wertig & 0.6 & 0.4 & 0.2 & 0 \\
\cline{2-6}
 & inakzep\-tabel & 0 & 0 & 0 & 0 \\
\hline
\hline
\end{tabulary}
\renewcommand{\arraystretch}{1}
\end{center}


Bringt man dies zusammen, ergibt sich folgender erster Eindruck:

\begin{center}
\renewcommand{\arraystretch}{1.5}
\begin{tabulary}{14cm}{|C|C C C C C C C C|C|}
\hline
 VON & V & \ra & VIA &  & & & & \ra & NACH \\
 \hline
 Frontend & $\bigstar$ & \ra & VIA &  & & & & \ra & NACH \\
\hline
\hline
\hline
 EasyABC  & \ra & ly2abc &  & & & & \ra &  \LaTeX\ \& ABC \\
 .xml & \ra & xml2abc mxml2abc & & &  &  & \ra & .abc \\
\hline
 .abc & \ra & abc2ly & & &  &  & \ra & .ly \\
 .xml & \ra & xml2ly musicxml2ly & & &  &  & \ra & .ly \\
\hline
 .xml & \ra & xml2pmx & & &  &  & \ra & .pmx \\
\hline
 .abc & \ra & abc2mtex & & &  &  & \ra & .tex \\
 .ly  & \ra & ly2abc & \ra &  .abc & \ra & abc2mtex & \ra & .tex \\
 .pmx & \ra & pmxab & & &  &  & \ra & .tex \\
 .xml & \ra & { xml2abc, mxml2abc} & \ra & .abc & \ra & abc2mtex & \ra & .tex \\
\hline
\end{tabulary}
\renewcommand{\arraystretch}{1}
\end{center}

\begin{small}
\begin{longtable}{|c|c|c|c|c|c||c||c||}

% was soll am Anfang oben scheinen:
\hline
  \multicolumn{2}{|c}{VON} & \multicolumn{2}{|c}{ÜBER} & 
  \multicolumn{2}{|c||}{NACH} & \multicolumn{2}{c||}{WERTUNG} \\
\hline
  Frontend & $\bigstar$ & Konvertierung & $\Sigma$  & \LaTeX\ \& & $\bigstar$ & 
  Güte\footnote{Die Güte wird in den folgenden Kapiteln empirisch ermittelt.} & 
  Rang\footnote{Der Rang ergibt sich über (stars(FE)+stars(BE)-num(Schritte))*Güte.}\\
\hline
\hline
\endfirsthead

% was soll bei jedem Seitenumnruch oben scheinen:
\multicolumn{4}{c}%
{\tablename\ \thetable\ -- \textit{Continued from previous page}} \\
\hline
\multicolumn{2}{|c}{VON} & \multicolumn{2}{|c}{ÜBER} & \multicolumn{2}{|c||}{NACH} & \multicolumn{2}{c||}{WERTUNG} \\
\hline
  Frontend & $\bigstar$ & Konvertierung & $\Sigma$  & \LaTeX\ \& & $\bigstar$ & 
  Güte\footnote{Die Güte wird in den folgenden Kapiteln empirisch ermittelt.} & 
  Rang\footnote{Der Rang ergibt sich über (stars(FE)+stars(BE)-num(Schritte))*Güte.}\\
\hline
\hline
\endfirsthead

% was soll bei jedem Seitenumnruch unter scheinen:
\hline \multicolumn{4}{r}{\textit{Continued on next page}} \\
\endfoot

% was am Ende erscheinen:
\hline
\endlastfoot

 Easy\-ABC & 4 & \ra\ & 0 &  ABC & 3 & 1 & (4+3-0)*1 = 7 \\
\hline
 Easy\-ABC & 4 & \ra\ abc2ly \ra\ ly2abc \ra\ & 2 &  ABC & 3 & 0.5 &  1\\
\hline
 Rose\-garden & 3 & \ra\ xml2abc \ra\ & 1 & ABC & 3 & 0.5 & 1,5\\
\hline 
Rose\-garden & 3 & \ra\ mxml2ly \ra\ ly2abc \ra\ & 2 &  ABC & 3 & 0.5 &  1\\
\hline
Rose\-garden & 3 & \ra xml2ly \ra\ ly2abc \ra\ & 2 &  ABC & 3 & 0.5 &  1\\
\hline
 Muse\-Score & 5 & \ra\ xml2abc \ra\ & 1 & ABC & 3 & 0.5 & 1,5\\
\hline
 Muse\-Score & 5 & \ra\ mxml2ly \ra\ ly2abc \ra\ & 2 &  ABC & 3 & 0.5 &  1\\
\hline
 Muse\-Score & 5 & \ra xml2ly \ra\ ly2abc \ra\ & 2 &  ABC & 3 & 0.5 &  1\\
\hline
 Denemo & 3 & \ra\ xml2abc \ra\ & 1 & ABC & 3 & 0.5 & 1,5\\
\hline
 Denemo & 3 & \ra\ mxml2ly \ra\ ly2abc \ra\ & 2 &  ABC & 3 & 0.5 &  1\\
\hline
 Denemo & 3 & \ra xml2ly \ra\ ly2abc \ra\ & 2 &  ABC & 3 & 0.5 &  1\\
\hline
 Canorus & 3 & \ra\ xml2abc \ra\ & 1 & ABC & 3 & 0.5 & 1,5\\
\hline
 Canorus & 3 & \ra\ mxml2ly \ra\ ly2abc \ra\ & 2 &  ABC & 3 & 0.5 &  1\\
\hline
 Canorus & 3 & \ra xml2ly \ra\ ly2abc \ra\ & 2 &  ABC & 3 & 0.5 &  1\\

\end{longtable}
\end{small}


% TODO nicht gegutachtete Wege: Fresco/ etc. nach Easyabc, da deaysy abc auch ein Texteditor


% this is only inserted to eject fault messages in texlipse
%\bibliography{../bib/literature}
